%% ============================================================
%%  Section: Robustness, Mechanisms, and Real Effects
%%  (2026-06-12, post-JMCB revision)
%%
%%  Organizes the new analyses per revision-plan-20260612.md:
%%    §1 2023 stress: flexible characteristics-by-quarter controls (RA2/RB1)
%%    §2 cohort timing + not-yet-adopted placebo (RB1)
%%    §3 extended pre-trends incl. COVID quarters (RA1/RA6)
%%    §4 residualized treatment (RA1)
%%    §5 alternative treatment scalings (RA5)
%%    §6 nonlinearity, symmetric tails, outlier (RB-m3/m4, RB-m1)
%%    §7 deposit destination: brokered/reciprocal + SOD spillovers (RB2)
%%    §8 lending and the bank-response channel (RA8, RB-m2)
%%
%%  Cross-references created:
%%    sec:robustness, sec:robustness:stress, sec:robustness:timing,
%%    sec:robustness:pretrends, sec:robustness:resid,
%%    sec:robustness:scaling, sec:robustness:nonlinearity,
%%    sec:robustness:destination, sec:robustness:lending
%% ============================================================

\section{Robustness, Mechanisms, and Real Effects}
\label{sec:robustness}

This section subjects the main results to the identification concerns raised
in Section~\ref{sec:identification} and develops the mechanism.  We proceed
from the most demanding threat---the 2023 banking stress---through the
specification of the treatment itself, and then turn to two questions any
deposit-discipline result must answer: where the deposits go, and whether the
funding pressure reaches the asset side of the balance sheet.

\subsection{The 2023 Stress: Characteristics-by-Quarter Controls}
\label{sec:robustness:stress}

If high-CECL banks were differentially exposed to the March 2023 stress
through observable balance-sheet characteristics, quarter fixed effects alone
would not remove the confound.  Table~\ref{tab:flexcontrols} therefore
re-estimates the baseline allowing progressively richer sets of 2022Q4 bank
characteristics to enter interacted with calendar-quarter fixed effects, so
that each characteristic traces its own unrestricted time path through the
tightening cycle and the stress episode.  Column~(1) repeats the baseline.
Column~(2) adds size and capital interactions; the estimate is essentially
unchanged ($-0.069$, $p < 0.10$ continuous; $-0.404$, $p < 0.05$ binary).
Column~(3) adds the deposit-mix characteristics most implicated in the 2023
runs---the uninsured share of deposits and the brokered share---and
column~(4) adds the full eight-characteristic vector including NPL, ROA, CRE
concentration, and unrealized securities losses.  Here the level estimates
attenuate toward zero and lose significance (continuous: $-0.043$ and
$-0.037$; binary: $-0.275$ and $-0.238$).  The event-study version under the
full characteristic set (Figure~\ref{fig:es_flexcontrols}) shows the same
post-adoption shape at compressed magnitude.

We report this attenuation prominently rather than bury it, because its
interpretation matters.  The controls in columns~(3)--(4) include the 2022Q4
\emph{uninsured deposit share}---a variable mechanically entangled with the
outcome.  Banks whose uninsured funding was about to be repriced by the
disclosure are banks with uninsured funding; allowing that characteristic an
unrestricted post-2023 time path absorbs not only any stress exposure it
proxies for but also part of the treatment effect itself.  This is the
classic bad-control problem, and it is why the composition design matters.
Table~\ref{tab:flexcontrols_td} applies the same characteristics-by-quarter
saturation \emph{within} the triple-difference, allowing the full 2022Q4
characteristic vector to shift the uninsured--insured gap flexibly in every
quarter.  The triple interaction survives essentially intact: $-0.136$ per
percentage point of CECL exposure ($p < 0.05$) and $-1.05$ percentage points
for high-CECL banks ($p < 0.01$), against baseline estimates of $-0.177$ and
$-1.25$.  Whatever differential stress exposure the 2022Q4 characteristics
carry, it does not explain why uninsured deposits fell \emph{relative to
insured deposits within the same bank} in proportion to the disclosed
adjustment.  A fully saturated triple-difference that adds all lower-order
interactions (Table~\ref{tab:triple_diff_saturated}) delivers identical
triple coefficients, confirming that the parsimonious specification does not
hide anything in the absorbed terms.

\subsection{Cohort Timing and a Placebo on Not-Yet-Adopted Banks}
\label{sec:robustness:timing}

The staggered rollout offers a second lever on the stress confound: 323
banks adopted in quarters other than 2023Q1, experiencing the March 2023
shock at event times other than $k = 0$.
Figure~\ref{fig:es_cohorts} overlays event studies estimated separately for
the early (2022Q3--Q4), main (2023Q1), and late (2023Q2--Q4) cohorts in
event time, and Figure~\ref{fig:es_cohorts_calendar} re-estimates them in
calendar time.  The divergence in each cohort tracks its own adoption
quarter rather than lining up on March 2023, which is difficult to reconcile
with a pure stress story.  The formal test has limited power and we present
it as such: dropping the entire 2023Q1 cohort leaves only 13 high-CECL banks,
and the resulting difference-in-differences (Table~\ref{tab:cohort_timing})
carries the right sign (continuous $-0.081$; binary $-0.133$) but is
imprecise.

Table~\ref{tab:placebo_notyet} reports a placebo we consider informative
precisely because it does not come out clean.  For the 231 late-cohort banks
that had not yet adopted as of 2023Q1, we regress the 2023Q1 change in
uninsured time deposits on the size of their \emph{future, still
undisclosed} Day-One adjustment.  If the adjustment were pure private
information until disclosure and the March stress orthogonal to it, the
coefficient would be zero.  Instead it is negative ($-0.101$, $p < 0.05$
continuous; $-0.574$, $p < 0.01$ binary), indicating that late adopters with
larger future adjustments already lost uninsured deposits in the stress
quarter.  Two readings are consistent with this pattern: depositors at these
banks responded in March 2023 to the same credit-risk fundamentals that CECL
subsequently capitalized into a single number, or some forward-looking
depositors anticipated the disclosure itself.  Either way, the placebo says
the level difference-in-differences around 2023Q1 cannot by itself separate
disclosure from stress---which is exactly why we rest the identification on
the characteristics-by-quarter controls and the within-bank composition
contrast of Section~\ref{sec:robustness:stress}, neither of which this
correlation can contaminate, and on the cohort-timing evidence above.

\subsection{Extended Pre-Trends, Including the Pandemic Quarters}
\label{sec:robustness:pretrends}

The baseline event window begins in 2020Q3.  Figure~\ref{fig:es_long}
extends the event studies back to 2018Q1, both including and excluding the
COVID-19 quarters (2020Q1--Q2), for uninsured and insured time deposits.
The point estimates over the five-year pre-period oscillate around zero with
no economically meaningful slope, and the post-adoption break remains the
dominant feature of every panel.  Formal joint tests of the pre-adoption
coefficients, however, reject marginally (uninsured: $F = 1.60$, $p = 0.04$
including COVID, $p = 0.04$ excluding; insured: $p = 0.07$ and $p = 0.05$),
driven by scattered individually significant leads deep in the pre-period
rather than by a trend toward the adoption-date divergence.  We report these
$p$-values rather than claim pristine parallel trends over a five-year
horizon for 4,300 banks; with this many estimated leads and tight standard
errors, statistical rejections of exact parallelism are close to inevitable,
and the economic content of the figures is the absence of any pre-existing
drift that could generate the sharp, persistent, correctly-timed
post-adoption divergence.  The COVID quarters themselves contribute noise
but no differential trend, addressing the concern that their exclusion from
the baseline window conceals unfavorable dynamics.

\subsection{Residualized Treatment}
\label{sec:robustness:resid}

To isolate variation in the Day-One adjustment orthogonal to everything a
depositor could have observed in 2022, we residualize \texttt{CECL Adj} on
the full expanded predictor set of
Table~\ref{tab:cecl_predictors_cross_section}, column~(2)---fundamentals,
deposit composition, and performance trends---and re-estimate the baseline
with the residual as treatment (Table~\ref{tab:resid_cecl};
Figure~\ref{fig:es_resid} reports the corresponding event study).  The point
estimates retain their signs but attenuate and lose significance (uninsured
deposits: $-0.052$ against a baseline of $-0.072$; interest expense and ROA
attenuate similarly).  We interpret this exercise as a bound rather than a
failure.  The depositor response \emph{should} load on the credit-risk
content of the disclosure, and the residualization strips out precisely the
component of the adjustment that is correlated with observable credit
quality---NPL ratios, deteriorating profitability trends---which is signal,
not confound.  What survives in the residual is the component of the
disclosure unpredictable even from a full 2022Q4 observable set; that its
effect is of the same sign and two-thirds the magnitude, though imprecise,
suggests the baseline is not driven purely by selection on observables,
while honestly conceding that the orthogonalized variation alone is too
coarse to re-establish significance.

\subsection{Alternative Treatment Scalings}
\label{sec:robustness:scaling}

Scaling the Day-One charge by equity could mechanically link treatment to
capitalization.  Tables~\ref{tab:alt_scaling_assets}
and~\ref{tab:alt_scaling_loans} re-estimate the main outcomes with the
charge scaled by total assets and by total loans, with binary indicators cut
at the same percentile as the baseline threshold.  The continuous estimates
preserve the baseline pattern: uninsured time deposits decline (assets
scaling: $-0.714$, $p < 0.10$; loans scaling: $-0.438$, $p < 0.10$, per
percentage point of the rescaled charge), interest expense rises, and ROA
falls, with the profitability effects significant at the 1~percent level
under both scalings.  The binary indicators are noisier, as expected when an
extreme-tail cutoff is applied to a differently shaped distribution.  The
equity denominator is thus not generating the results, though we retain it
as the baseline because equity scaling measures what depositors plausibly
care about: the fraction of the loss-absorbing buffer consumed by the newly
recognized expected losses.

\subsection{Nonlinearity, Symmetric Tails, and the Outlier}
\label{sec:robustness:nonlinearity}

The treatment distribution has a mass point at zero---roughly three-fifths
of banks recorded no Day-One adjustment---so we examine the dose-response
with economically defined bins rather than degenerate quantile cuts.
Figure~\ref{fig:cecl_bins} plots post-adoption coefficients for uninsured
time deposits by adjustment category, with exactly-zero banks as the
reference group: negative (favorable) adjustments, small positive (below the
median positive adjustment, 0.85 percent of equity), large positive (median
to 90th percentile, 3.2 percent), and extreme positive (above).  The
response is monotone in the adverse direction: $-0.24$ for small positive
($p = 0.06$), $-0.40$ for large positive ($p < 0.01$), and $-0.56$ for
extreme positive adjustments ($p < 0.05$).  Alternative binary cutoffs at
the 75th and 90th percentiles deliver significant estimates of $-0.28$ and
$-0.29$ (Table~\ref{tab:alt_cutoffs}), so the baseline tail threshold is not
doing the work.

The left tail is more interesting and we do not oversell it.  If depositors
read the adjustment as a symmetric signal, banks with \emph{negative}
(favorable) Day-One adjustments should attract uninsured inflows.  They do
not: the negative-adjustment bin carries a coefficient of $-0.27$
($p < 0.05$), and a bottom-decile indicator (most favorable adjustments,
top decile excluded for a clean contrast) yields $-0.27$ ($p < 0.10$,
Table~\ref{tab:alt_cutoffs}, column~4).  The response to the disclosure is
asymmetric: adverse news drives outflows that scale with its size, while
favorable news generates no offsetting inflow---if anything, any sizable
adjustment in either direction associates with mild outflow relative to
banks reporting exactly zero.  The asymmetry is consistent with
loss-focused depositor attention---an account in which the disclosure
triggers scrutiny rather than a two-sided Bayesian update---and parallels
the asymmetric reaction to good and bad bank news documented elsewhere in
the disclosure literature.  It also cautions against linear extrapolation
of the continuous estimates into the left tail.

Finally, Figure~\ref{fig:cecl_equity_distribution} shows a single visible
outlier in the right tail.  Table~\ref{tab:outlier} drops the bank with the
largest adjustment: the continuous estimate moves from $-0.072$ to $-0.068$
($p < 0.10$) and the binary estimate from $-0.394$ to $-0.381$
($p < 0.05$).  No single observation drives the results.

\subsection{Where Do the Deposits Go?}
\label{sec:robustness:destination}

A composition shift from uninsured to insured deposits raises the question
of mechanism: are depositors leaving, or re-contracting into insured form?
Table~\ref{tab:brokered} examines the wholesale-funding margin directly.
After adoption, high-CECL banks increase insured brokered deposits by 0.52
percentage points of assets ($p < 0.05$) and reciprocal deposits---the
market arrangement that purchases full insurance coverage for large
accounts---by 0.089 percentage points per unit of continuous CECL exposure
($p < 0.10$), while \emph{uninsured} brokered funding does not move (point
estimates indistinguishable from zero).  The event studies
(Figures~\ref{fig:es_brokered_ins} and~\ref{fig:es_reciprocal}) show flat
pre-trends and post-adoption accumulation.  High-CECL banks did not lose
funding so much as repurchase it in insured form, at the higher rates
documented in Section~\ref{sec:results:financial}.  This is the
market-based insurance margin emphasized by \citet{kim2024insurance}, and
it implies that part of the measured insured-deposit increase reflects
managed substitution by banks rather than spontaneous reallocation by retail
depositors.

To gauge how much of the composition result this substitution accounts for,
Table~\ref{tab:td_excl_brokered} re-estimates the triple-difference
excluding the quartile of banks with the largest post-adoption growth in
insured brokered deposits.  The triple interaction shrinks toward zero and
loses significance ($-0.063$ continuous, $-0.46$ binary), indicating that
the within-bank composition shift is concentrated precisely among the banks
that actively replaced uninsured funding with insured wholesale money.  Far
from undermining the mechanism, this localization completes it: the
disclosure repriced uninsured funding, and the banks facing the largest
repricing are the ones that re-intermediated through the insured wholesale
market.

The remaining piece is the local market.  If withdrawn uninsured deposits
re-intermediate nearby, treated-control spillovers would violate the
stable-unit assumption and bias the difference-in-differences \emph{toward}
zero, since control banks partially absorb the treatment.
Table~\ref{tab:spillover} tests this with branch-level Summary of Deposits
data.  Among low-CECL banks, annual deposit growth after 2023 increases
with the bank's 2022 leave-out exposure to high-CECL competitors in its
counties ($+2.4$ percentage points per unit of exposure, $p < 0.10$); at
the county level, deposits at low-CECL banks' branches grow faster after
2023 in counties with larger high-CECL deposit shares ($+5.2$, $p < 0.01$).
Both estimates use annual June-to-June data and modest cross-sectional
variation, so we read them as qualitative confirmation: deposits leaving
high-CECL banks partly land at local competitors, the direction of
spillover that makes our within-market estimates conservative.

\subsection{Lending and the Bank-Response Channel}
\label{sec:robustness:lending}

A growing literature documents that banks respond to CECL itself---through
provisioning, capital management, and credit supply
\citep{granja2025current, chen2022does, kim2026current}---which raises an
alternative reading of our results: perhaps high-CECL banks changed their own
behavior at adoption, and deposit movements respond to the bank's actions
rather than to depositors' reading of the disclosure.  Three observations
bound this channel.  First, the timing works against it: deposit outflows
appear in the first post-adoption quarters, simultaneous with the first
possible strategic response and before any altered lending policy could
plausibly be observed by depositors.  Second, the flat pre-trends rule out
the version in which banks adjusted ahead of adoption.  Third, and most
directly, the funding pressure \emph{precedes} the lending response in the
event-time data: Table~\ref{tab:lending} shows that quarterly total loan
growth at high-CECL banks slows by 0.65 percentage points after adoption
($p < 0.01$), concentrated in commercial real estate ($-0.90$, $p < 0.01$),
with commercial and industrial lending negative but imprecise; the
event study (Figure~\ref{fig:es_loan_growth}) shows the slowdown building
over the post-period rather than jumping at adoption, as one would expect if
lending adjusts to a funding shock with a lag.  These real effects close the
loop of the paper: a mandatory disclosure repriced uninsured funding, banks
replaced it with more expensive insured money, and the higher marginal cost
of funds passed through to credit supply, the margin emphasized by
\citet{granja2025current}.
